% Одностраничный отчёт по образцу САО.
% Все поля, которые обычно приходится менять перед печатью, собраны здесь.
\documentclass[11pt,letterpaper]{article}

\usepackage[T2A]{fontenc}
\usepackage[utf8]{inputenc}
\usepackage[russian]{babel}
\usepackage{amsmath}
\usepackage{geometry}
\usepackage{microtype}
\usepackage{tabularx}

\geometry{
  top=22mm,
  bottom=15mm,
  left=20mm,
  right=20mm
}

% ---------------------------------------------------------------------------
% Редактируемые поля
% ---------------------------------------------------------------------------
\newcommand{\StudentFullName}{Буканов Иван Алексеевич}
\newcommand{\StudentShortName}{Буканов И. А.}
\newcommand{\StudentAffiliation}{ФКИ МГУ}
\newcommand{\PracticeStart}{08 июля}
\newcommand{\PracticeEnd}{24 июля 2026 г.}
% Дата в блоке подписей: можно написать здесь любую строку целиком.
\newcommand{\ReportDate}{«24» июля 2026 г.}
\newcommand{\SupervisorFullName}{Макаров Д. И.}
\newcommand{\ProjectTitle}{Калибровка метода флуктуаций поверхностной яркости по данным JWST/NIRCam}

\setlength{\parindent}{0pt}
\setlength{\parskip}{3pt}
\setlength{\abovedisplayskip}{4pt}
\setlength{\belowdisplayskip}{4pt}
\setlength{\emergencystretch}{2em}
\pagestyle{plain}

\begin{document}

\begin{center}
  {\LARGE Отчёт о работе над научным проектом}
\end{center}

\vspace{6mm}

Студент \StudentFullName, \StudentAffiliation

\vspace{2mm}

Я прошёл практику в Лаборатории внегалактической астрофизики и космологии
САО РАН в период с \PracticeStart{} по \PracticeEnd{} под руководством
заведующего лабораторией \SupervisorFullName{} по теме
«\ProjectTitle».

В ходе работы я разработал на языке Python воспроизводимый алгоритм измерения
расстояний до галактик методом флуктуаций поверхностной яркости
(Surface Brightness Fluctuations, SBF). Были обработаны публичные изображения
14 галактик ранних морфологических типов из программы JWST GO-3055. Данные
NIRCam в фильтре F150W использовались для измерения SBF-сигнала, а изображения
F090W --- для независимой проверки возможной зависимости абсолютной
флуктуационной звёздной величины от цвета звёздного населения.

Для каждой галактики были выполнены оценка и вычитание фона, поиск и
маскирование компактных источников, построение гладкой модели по системе
эллиптических изофот и получение остаточного изображения. Измерения проводились
в двух круговых кольцах с радиусами $8.2$--$16.4$ и $16.4$--$32.8$ угл. сек,
что соответствует геометрии опубликованного анализа программы GO-3055.
Амплитуда флуктуаций определялась в пространстве пространственных частот
путём аппроксимации спектра мощности моделью
\[
  P(k)=P_0E(k)+P_1,
\]
где $E(k)$ учитывает функцию рассеяния точки и маску рабочей области,
$P_0$ содержит мощность SBF и неразрешённых компактных источников, а $P_1$
описывает белый шум. Из $P_0$ вычиталась оценка вклада неразрешённых шаровых
скоплений и фоновых галактик $P_r$.

Полученные флуктуационные величины были откалиброваны по индивидуальным
расстояниям, измеренным методом вершины ветви красных гигантов (TRGB) в работе
TRGB--SBF Project IV. Учтено поглощение в Млечном Пути; исследованы ошибки,
связанные с фоном, PSF, маскированием, вкладом $P_r$ и крупномасштабной
структурой остатков. Для полной выборки получен цветонезависимый нуль-пункт
$\overline{M}_{150,0}=-3.172$ зв. вел. с внутренним разбросом $0.084$ зв. вел.
Перекрёстная проверка с исключением по одной галактике дала RMS $0.0955$
зв. вел., что соответствует разбросу около $4.4\%$ по расстоянию, при
практически нулевом среднем смещении относительно TRGB.

Линейная зависимость от цвета была проверена отдельно, однако не улучшила ни
информационный критерий, ни точность перекрёстной проверки. Поэтому основным
результатом принята калибровка с постоянным нуль-пунктом. Для контроля
устойчивости дополнительно выполнены тесты нормировки PSF, чувствительности к
$P_r$, синтетического восстановления известной SBF-амплитуды и качества
остаточных изображений. Подготовлены итоговые таблицы, диагностические графики,
пакетный сценарий обработки и черновик научной статьи. Дальнейшая работа будет
связана с улучшением двумерной модели галактик и применением метода к более
широкой выборке JWST, включая программу J-Virgo GO-7763.

\vfill

\noindent
\begin{tabularx}{\textwidth}{@{}X X@{}}
  \rule{38mm}{0.4pt} & Руководитель проекта: \\
  \ReportDate & \rule{38mm}{0.4pt}\,/ \SupervisorFullName \\
  \\[3mm]
  \rule{38mm}{0.4pt} & Исполнитель: \\
  \ReportDate & \rule{38mm}{0.4pt}\,/ \StudentShortName
\end{tabularx}

\end{document}
